\subsection{Sylvesterの判定法1}

\begin{definition}
  \label{def:Matrix.leadingSubmatrix}
  \lean{Matrix.leadingSubmatrix}
  \leanok
  $M_k$と書き，$M$の$k$次主小行列を表す．
\end{definition}

$X$を任意の添字集合で$|X| = m$，$e : \map{X}{I}$を写像とする．
行列$M_{e, e} \in \M_m(\R)$を
\begin{equation}
  (M_{e, e})_{ij} := M_{e(i),\ e(j)}
\end{equation}
となるように定める．

\begin{lemma}
  \label{lem:IsHermitian.submatrix_injective}
  \lean{IsHermitian.submatrix_injective}
  \leanok
  エルミート行列$M$に対し，$M_{e, e}$もまたエルミートである．
\end{lemma}

\begin{lemma}
  \label{lem:PosDef.submatrix_injective}
  \lean{PosDef.submatrix_injective}
  \leanok
  \uses{lem:IsHermitian.submatrix_injective}
  写像$e$は単射であるとする．
  正定値行列$M$に対し，$M_{e, e}$もまた正定値である．
  つまり，正定値行列の部分行列もまた正定値である．
\end{lemma}

\begin{proof}
  部分行列がエルミートであることは上の補題~\ref{lem:IsHermitian.submatrix_injective}から従う．

  $ 0 \ne x := \transpose{(x_1, \ldots, x_m)} \in \R^m$とする．
  \begin{equation}
    y_i = \begin{cases}
      x_j, & e(j) = i,\\
      0, & \mathrm{otherwise}
    \end{cases}
  \end{equation}
  とし，ベクトル$y \in \R^n$を$y := \transpose{(y_1, \ldots, y_n)}$と定める．
  すると，それぞれの二次形式は一致する：
  \begin{equation}
    x^* M_{e, e} x = y^* M y. \label{eq:quad_form_xy}
  \end{equation}
  $y = 0$ならば$x = 0$なので，対偶をとれば$y \ne 0$がわかる．
  $M$の正定値性より，式~\eqref{eq:quad_form_xy}は正であるから，$M_{e,e}$の正定値性が示せた．
\end{proof}

Sylvesterの判定法の$\ref{enu:Sylvester_posdef} \implies \ref{enu:Sylvester_det}$を示す：

\begin{theorem}
  \label{thm:posDef_leading_minors_pos}
  \lean{posDef_leading_minors_pos}
  \leanok
  \uses{lem:PosDef.submatrix_injective}
  $M$が正定値行列であるとき，任意の$k \le n$に対し，$\det M_k > 0$である．
\end{theorem}

\begin{proof}
  正定値行列の行列式は正であるので，$M_k$が正定値であることを言えば十分である．
  補題~\ref{lem:PosDef.submatrix_injective}より$M_k$は正定値であるから，定理の主張が従う．
\end{proof}